%!Mode:: "TeX:UTF-8"
% !TeX program = xelatex

\documentclass{setp-new}
% Ctex 可以使用pdflatex或者xelatex编译
%此模板需调用宏包cjoe-new.cls, 请确保此宏包与tex文件在同一文件夹下
\usepackage{float}
\usepackage{pifont}
\usepackage{mathtools}
\usepackage{hhline}

\jolname{系统工程理论与实践}{Systems Engineering --- Theory \& Practice}

\newcommand{\sz}{\small}


\begin{document}

\Volume{xxxx}{0}{0}{0}  % 年、月、卷、期
\PageNum{1}              % 起始页码
\PaperDOI{10.12011/SETPXXXX-XXXX}  % DOI,编辑部给出
\Category{XXX{\bf 请作者给出}}{A}


\TitleMark{作者A, 等: 系统工程理论与实践论文\LaTeX 模板说明}  % 页眉, 如只有一位作者,请去掉``等" 若题目过长, 请去掉副标题

\Title{系统工程理论与实践论文\LaTeX 模板说明}%标题

\Author{作者A{\fontsize{12pt}{\baselineskip}\selectfont $^{1}$}, 作者B{\fontsize{12pt}{\baselineskip}\selectfont $^{2}$}, 作者C{\fontsize{12pt}{\baselineskip}\selectfont $^{1}$}, 作者D{\fontsize{12pt}{\baselineskip}\selectfont $^{2}$}}%姓名为两个字用 \quad 隔开
{(1.学校名称~ 所在学院, 所在城市~ 邮编; 2. 中国科学院 ~数学与系统科学研究院, 北京  100190)}%作者及单位

\Abstract{本文给向《系统工程理论与实践》投稿或已录用的作者提供一个中文\LaTeX{}模版,详细说明了本刊编排要求,请作者认真阅读,并在相应的位置填入相应的内容既可, 模板版面设置参数不允许修改.}

\Keywords{关键词1; 关键词2; 关键词3}%关键词用英文分号隔开

\ETitle{The guide of the \LaTeX{} template for preparing the manuscript of Systems Engineering --- Theory \& Practice} %英文标题

\EAuthor{ZUO Zhe A$^1$, \ ZUO Zhe B$^2$, \ ZUO Zhe C$^1$,  \ ZUO Zhe D$^2$} % 作者姓名的拼音,姓在前,名在后,姓大写.
{(1. University, City~ Postal Code, Country;\\
2. Chinese Academy of
Sciences, Academy of Mathematics and Systems Science, Beijing
100190, China)}

\EAbstract{ This article provides a Chinese
\LaTeX{} template for authors who have submitted or accepted to ``Systems Engineering --- Theory \& Practice", detailing the layout requirements of this journal, please read carefully and fill in the corresponding content in the corresponding position. }

\EKeywords{key word1; key word2; key word3}%关键词用英文分号隔开, 首字母小写


% -------- \EditorNote 放到 title 的后面

\EditorNote{\textbf{收稿日期:}\ XXXX-XX-XX}  % 脚注,收稿日期 编辑部给出

\EditorNote{\textbf{作者简介:} 姓名(出生年月--), 职称, 学位, 学术称号, 研究方向, E-mail: XXX; 姓名(出生年月--), 职称, 学位, 学术称号, 研究方向, E-mail: XXX.} %所有作者的作者简介

\EditorNote{\textbf{基金项目:}\ 基金名称(项目编号); 国家自然科学基金(71774051)\\
\textbf{Foundation item:} Fund Name (XXXXXXXX); National Natural Science Foundation of China (71774051)}%中文基金相对应的英文翻译.

\CiteAuthor{作者A, 作者B, 作者C, 等.}{Author A, Author B, Author C, et al.}



\maketitle
\section{引言(必读)}
%一级标题

本模版是录用后作者自行初排的模板,也可投稿使用.

本模板所适用的编译环境是 TeXLive 2017以上的版本.

本模板应使用XeLaTeX 编译.

请优先使用本模板中的代码.



作者须知:


1) 排版过程中 请不要自行插入多余的任何符号如\verb|\par|, \verb|~| \verb|\,|等.


2) 全文标点符号除顿号外均用``英文"中的标点, 顿号用中文标点并且后面加个\verb|~|, 引号用`` ".

3) doi号、~收稿日期、~年、~月、~卷、~期、~起始页码均由编辑修改,作者无需修改,
中图分类号请作者在指定位置给出.

4)模板版面设置参数不允许修改.

5) 正文中的公式、~图、~表、~定理、~推理、~定义、~命题等, 按在全文中的先后顺序直接按顺序号编号, 如:  定理1, 定理2, $\cdots$$\cdots$, 不以章节编号.

6)文章若有附录, 请作者将附录直接附在正文的最后面, 文章发表时纸质期刊仅刊发正文部分; 本刊官网、~中国知网等电子版将刊发全文(正文+附录). 在不影响可读性情况下, 文章背景介绍、~问题描述、~分析过程等尽量言简意赅, 请尽可能将证明过程、~算法步骤(/步骤图)、~数据分析过程、~源数据、~成熟的概念或模型等放入附录.

7) tex文件中有些灰色字体的说明,可供阅读.

\section{《系统工程理论与实践》格式规范(参照)}


\subsection{文章标题}


1)中文文题　

不超过20个汉字, 不采用非通用的缩写词.

2)英文文题　

不超过10个实词, 不采用非通用的缩写词.

\subsection{文章作者}


1) 作者简介　

按模板填入既可.

2)作者单位
　　

a. 作者单位要具体到院系、~实验室或研究所;
b. 单位英文名称须是规范、~标准的全称, 不采用缩写;
c. 注意检查中英文单位的名称、~邮编、~序号标注是否一致.


\subsection{基金项目}


1)基金项目须用国家统一的标准名称, 勿用简化形式表述.

2)基金项目后须附项目编号.

3)基金项目须全部译成英文.

\subsection{摘要}


1)中文摘要

a. 摘要应具有独立性和自明性, 重点叙述文章创新之处, 避免描述常识性内容, 基本要素包括研究的目的、~方法、~结果、~结论;
b. 摘要中有英文缩略语的, 应在第一次出现时给出英文全称;
c. 字数在200$\sim$300字为宜.
　　

2)英文摘要
　　　

a. 英文摘要须与中文摘要内容一致, 英文请仔细检查, 做到语法正确、~文字流畅;
b. 摘要中出现的英文缩写要用全称, 缩写放在全称后面, 用括号括起.

\subsection{关键词}

关键词数量以3$\sim$8个为宜, 按涵盖范围从大到小排列, 关键词为常用的检索词, 以便于数据库检索.

关键词用全称, 不单独采用缩写, 缩写须放在全称后面, 以括号括起.

\textbf{注: 2.1$\sim $2.5节应保持中英文一致}



\subsection{中图分类号}
中图分类号从《中国图书馆分类法》(第五版)中查得.



\subsection{脚注}

 脚注可使用\verb|\footnote{  .}|.
 %示例: 系统工程理论与实践\footnote{系统工程理论与实践论文\LaTeX 模板说明.}.注意上标1要在标点符号前.

\section{模板说明+示例}

\subsection{定理说明+示例}
%二级标题

\subsubsection{定理说明(必读)}
关于定理1、~引理1、~证明1、~ 假设1、~ 定义1、~ 例1、~ 实验1、~命题1、~算法1等都可运用代码 \textbf{定理1}\ \ 手动生成,注意后面均不用加\verb|:|.

\subsubsection{示例}%三级标题

\textbf{证明}\ \ $\cdots\cdots$;
\textbf{假设2}\ \  $\cdots\cdots$

\subsection{公式说明+示例}

\subsubsection{公式说明(必读)}

文章中的公式请按照例子中给出的代码编排, 若例子中不包含, 可自行编排, 建议尽可能使用简短的代码编排.


1)公式太长需换行的, 前一行的行末必须是运算符, 然后方可换到下一行;

2)公式序号须全文连贯顺排, 不可出现2(a), a1, A1等形式, 必须按照(1),  (2), $\cdots$$\cdots$编排;

3) 每个公式后面都应该有标点符号, 一般``其中"前面的公式用逗号, ``其中"应顶格排列,不要空两格.  其余用句号. 多行公式应用逗号, 最后一排用句号.

{\bf 重点}: 4) 正文中出现的所有计算量, 必须使用正确的字体: 凡表示变量或一般函数的外文字母, 用斜体. 如: $y(k)=c(k)+e(k)$, 多字母变量建议用正体, 如sym, st等;斜体(Latex排版中使用命令\verb|$ $|, 不要用\verb|\emph{}|), 向量(矢量)和矩阵用粗斜体(Latex排版中使用命令\verb|\bm|), 常量用正体(Latex排版中使用命令\verb|\rm|).

{\bf 重点}: 5) 特殊函数的符号、~圆周率$\pi$、~自然对数底e、~微分符号d、~共轭符号H、~kW、~H$_2$O、~sin、~lim、~max、~sup、~diag、~log、~时间s, 长度m, 微分d, 指数e, 增量$\delta$、~转置符号T等均用正体.

6) 单位和数字间应空一格 如8 t、 ~29 kW等.


\subsubsection{示例}

单行公式 \verb|\begin{equation}\end{equation}|

多行公式 \verb|\begin{eqnarray}\end{eqnarray}|

矩阵\verb|\begin{dcases}\end{dcases}|

\noindent{\bf 种类1\ \ 公式}


1)单行公式

\ding{192} 带公式序号
 \begin{equation}
    M = mN^{1}_{R} + mN^{1}_{S}/\rho.
\end{equation}

\ding{193} 不带公式序号
 \begin{equation*}
    M = mN^{1}_{R} + mN^{1}_{S}/\rho.
\end{equation*}

或
 \begin{equation}
    M = mN^{1}_{R} + mN^{1}_{S}/\rho.\nonumber
\end{equation}

或
$$    M = mN^{1}_{R} + mN^{1}_{S}/\rho.$$



2) 多行公式(应左对齐)

\ding{192} 带公式序号
\begin{eqnarray}
&&\Re_{\rm e m p}(f)=\frac{1}{N} \sum_{i=1}^{N}
\ell\left(f\left(x_{i}; \theta\right), y_{i}\right),\\&&
f^{*}={\arg \min}_{f \in H}   \Re_{\rm e m p}(f).
\end{eqnarray}

\ding{193} 不带公式序号
\begin{eqnarray*}
&&\Re_{\rm e m p}(f)=\frac{1}{N} \sum_{i=1}^{N}
\ell\left(f\left(x_{i}; \theta\right), y_{i}\right),\\&&
f^{*}={\arg \min}_{f \in H}   \Re_{\rm e m p}(f).
\end{eqnarray*}%或者去掉*,加\nonumber 也可以

3)公式分行

\ding{192} 公式在符号后断开.
\begin{equation}
\begin{aligned}
\Pi_i = & \{ a - b[ ( 1 - \sigma_{1}
)q_{1} +( 1 -\sigma_{2} )q_{2} +
(\sigma _{1} +\sigma_{2} ) ]
\}[ ( 1 -\sigma_{i} )q_{i} + \sigma
_{i} ]( 1 - \theta  )-\\& ( 1 -
\sigma _{i} )(cq_{i} - W ) - W.
\end{aligned}\end{equation}

或
\begin{align}
\Pi_i = & \{ a - b[ ( 1 - \sigma_{1}
)q_{1} +( 1 -\sigma_{2} )q_{2} +
(\sigma _{1} +\sigma_{2} ) ]
\}[ ( 1 -\sigma_{i} )q_{i} + \sigma
_{i} ]( 1 - \theta  )-\\& ( 1 -
\sigma _{i} )(cq_{i} - W ) - W.
\end{align}

\ding{193} 公式内有分式括号括不住时应使用\verb|\left \right| 使括号变大, 但是行内公式不要加.
\begin{align}
g( \eta_{t}| I_{t-1} )=\frac{1}{\sqrt{2\pi }}{{\rm e}^{-\frac{\eta _{t}^{2}}{2}}}\left( 1+\frac{s_{t}^{*}}{3!}( \eta _{t}^{3}-3{\eta _{t}} ) \right).
\end{align}

\noindent{\bf 种类2\ \ 矩阵}


1) 矩阵中元素居中对齐;
\begin{align}
\left[\begin{array}{cccc}
1& 1 &\dots &1\\
1& 1&\dots &1\\
\vdots&\vdots& \ddots   &\vdots\\
1& 1&\dots &1
\end{array}\right].\end{align}


2) 两个公式左边对齐.
\begin{align}
&A_1=\left[\begin{array}{cc} 1& -2\\-2 &1
\end{array}\right],\\
&A_2=\left[\begin{array}{ccc} 5&-5&-5\\-5&5 &5
\end{array}\right].
\end{align}

3) 有分式时应注意加\verb|\displaystyle|使分式变大, 分式变大只适用于公式环境, 不要用在行内.
\begin{align}
&A_3=\left[\begin{array}{ccc}\vspace{2mm} \displaystyle\frac{1}{2}&-1&-1\\ \displaystyle\frac{1}{2} &1 &1
\end{array}\right].
\end{align}


\noindent{\bf 种类3\ \ 带有大括号的例子}
\begin{equation}\begin{dcases}
\Delta _{l1t} = y_{lit}^1 - y_{lit}^0,\\
\Delta _{h1t} = y_{hit}^1 - y_{hit}^0.
\end{dcases}\end{equation}


\subsection{插图+示例}

\subsubsection{插图(必读)}

文章中的插图请按照例子中给出的代码编排, 若例子中不包含, 可自行编排, 建议尽可能使用简短的代码编排.


1)每幅图都必须有图题, 图题须有惟一性, 不可以与其他图题相同, 一套图中有多个分图的, 须在每个分图的下方标明分图的图题.

2) 图的图例须是中文, 置于图下方、~图题上方的位置.

3) 图有纵横坐标的, 相应的纵横坐标须有量和单位, 例如: 时间/s.

4) 正文中第一次提到图后, 再出现图; 正文中引用图时要明确图号, 不能用``如下图"等方式引用.

{\bf 重点:} 5) 插图要求: (文中的图片请严格遵守以下要求, 若达不到要求后期还会让其修改)


\ding{192} 文中的图应为矢量图, 应以.pdf 的格式插入.矢量图标准:放大缩小不会影响清晰度.即放大后线条、~字母等不会变虚,打印机打印出来不会模糊;

矢量图生成方式:

Word、~Excel、~ Powerpoint、~ Visio做的图都可以直接另存为或打印为pdf发过来. 可同时将docx, xlsx, pptx, vsd文件发过来以便编辑部对图片进行格式处理. Matlab做的图请在Matlab中生成图时直接另存为pdf格式发过来.

\ding{193} 图片不能有底纹;

\ding{194} 本刊为黑白印刷, 线条、~图例等应以线型或颜色＋线型的方式区分, 不能只用颜色区分;

\ding{195} 图文中的表述应与文中一致, 包括字母的正斜体, 中英文表述; 图中中文: 宋体; 英文: Times New Roman (运算符号、~阿拉伯数字、~常量), Times New Roman  Italic (单字母变量); 图文中的所有文字大小比图题小一号;

\ding{196} 图中出现的变量、~矢量、~常量, 要与中文的表述一致, 须用相应的斜体、~黑斜体、~正体表示;

\ding{197} 不要使用Photshaop软件编辑加工, 推荐使用pdf软件自带的编辑工具或illustrator CS5编译.


{\bf 重点:} 6) 图注的代码例子已给出,不要再用其它的.

{\bf 重点:} 7) 图文、~图表并排的代码可参照双图并排的代码.



\subsubsection{示例}

1)单个图

\begin{figure}[H]
\centering
\includegraphics[scale=.6]{1.pdf}\\\footnotesize{注: $\cdots$$\cdots$.}
\caption{外部媒体}
\label{fig1}
\end{figure}


2)双图并排(不带小标题)

\begin{figure}[H]
\begin{minipage}{0.5\linewidth}
\centering
\includegraphics[scale=.5]{2a.pdf}
\caption{外部媒体}
\label{fig2}
\end{minipage}
\begin{minipage}{0.5\linewidth}

\centering
\includegraphics[scale=.5]{2b.pdf}
\caption{外部媒体}
\label{fig2}
\end{minipage}
\end{figure}


3)双图并排(带小标题)

\begin{figure}[H]
\begin{minipage}{0.5\linewidth}
\centering
\includegraphics[scale=.5]{2a.pdf}\\
{\sz (a)  }
\end{minipage}
\begin{minipage}{0.5\linewidth}

\centering
\includegraphics[scale=.5]{2b.pdf}\\
{\sz (b)  }
\end{minipage}
\caption{外部媒体}
\end{figure}



\subsection{表格说明+示例}


\subsection{表格说明(必读)}

文章中的插图请按照例子中给出的代码编排, 若例子中不包含, 可自行编排, 建议尽可能使用简短的代码编排.


1) 每个表必须有表题, 且各表题之间不重复; 表题在表格的正上方. 表中参量应标明单位.

2) 表中的量和单位必须齐全.

3) 表中文字必须是中文, 否则请译成中文.

4) 正文中提到表以后, 再出现表; 正文引用表时要明确表号, 不能用``如下表"等方式引用.


{\bf 重点:} \ding{192} 表使用三线表, 较复杂的表可增加辅助线,表第一列应左对齐.


 \ding{193} 表格里的公式不要加\verb|\left \right|、~ \verb|\displaystyle|、~\verb|\dfrac|等使其变大.

 \ding{194}
表格宽度过长时可用\verb|\tabcolsep 2pt或\resizebox{\textwidth}{!}{}|使其缩短.\\
表格高度过高时可用\verb|\renewcommand{\arraystretch}{0.98}|使其缩短. 注意只能特殊情况下使用.

\ding{195} 表格里的***、~**、~*应用例子给出代码,不要用\verb|\ast|.


5) 表中的线请按照例子给出的运用, 不要全部都是\verb|\hline|.

6)表注的代码, 例子已给出,请参照, 不要使用其他的.

7)不要把表截成图插进来, 应用\LaTeX 绘制表.

8)  一页排不下的表格, 可使用续表排在下一页.续表表题只需写表1\ \ (续),不需再列一遍表题.



\subsection{示例}

1)正常情况
\begin{table}[H]
\centering
\caption{到达涨跌停板的超出概率与股票交易的活跃程度}
\label{tab1}
\tabcolsep 9pt
\begin{tabular}{lcccc}%第一列左对齐, 其他列居中
\toprule
& \multicolumn{4}{c}{$\text{Reach\_Limit}_{i, t}$} \\
\cmidrule{2-5}
 & (1) & (2) & (3) & (4) \\
\midrule
\multirow{2}{*}{$\text{Theo\_Prob}_{i, t}$} & 1.500***& 1.600*** & 1.500***& 1.600*** \\
 & (0.027) & (0.028) & (0.028) & (0.028) \\
\multirow{2}{*}{$\text{Direction}_{i, t}$} & & $-$0.150***& & $-$0.150*** \\
 & & (0.013) & & (0.013) \\
\bottomrule
\end{tabular}
\tabnote[.63\linewidth]{注: *、~**、~***分别表示10\%、~5\%和1\%的显著性水平.}
\end{table}

2) 表格宽度过长时可用\verb|\tabcolsep 2pt| (手动调整,最低值为2pt, 若还是超出版心建议选择3)的排法)或
\verb|{\textwidth}{!}{}(自适应)|使其缩短.


表格高度过高时可用\verb|\renewcommand{\arraystretch}{0.98}| (不可低于0.95)使其缩短. 注意只能特殊情况下使用(即表格只需要缩小一点间距就可放在当前页面时).


\begin{table}[H]
\centering
\caption{应急代储情形下运营商和批发商的均衡决策}%\resizebox{\textwidth}{!}{
\tabcolsep 3pt
\renewcommand{\arraystretch}{0.98}
\begin{tabular}{lcccc} \toprule
     \multirow{2}{*}{条件} & \multirow{2}{*}{单位交易费用$(\alpha^{E})$} &  \multirow{2}{*}{储备数量$(Q_{s}^{E})$} & \multicolumn{2}{c}{批发商$i$的订货数量}  \\\cmidrule(r){4-5}
     &&& 正常情形$(q_{iG}^{E})$ & 突发事件情形$(q_{iL}^{E})$ \\

\hline
  $0<t\leq t_{1}$ & $\frac{1-(1+k\rho-\rho)w_{G}}{2}$ & 0 &$ \frac{1+(-1+k\rho-\rho)w_{G}}{2(m+1)b_{w}} $ &  $\frac{1+(1+k\rho-\rho-2k)w_{G}}{2(m+1)b_{w}}$ \\
 $t_{1}<t<t_{2}$  & $\frac{1-(1+k\rho-\rho)w_{G}}{2}$ & $ \frac{(mw_{G}k+1)\rho-(m+1)(\bar{c}-t)}{2 b_{w}\rho} $  &$ \frac{1+(-1+k\rho-\rho)w_{G}}{2(m+1)b_{w}} $ & $\frac{1+(1+k\rho-\rho-2k)w_{G}}{2(m+1)b_{w}}-\frac{Q_{s}^{E}}{m+1}$  \\
 $t_{2}\leq t<\bar{c}$ & $\frac{1-w_{G}}{2}$ & $\frac{\rho+t-\bar{c}}{2 b_{w}\rho}$ &$ \frac{1-w_{G}}{2(m+1)b_{w}} $ & 0  \\  $0<t\leq t_{1}$ & $\frac{1-(1+k\rho-\rho)w_{G}}{2}$ & 0 &$ \frac{1+(-1+k\rho-\rho)w_{G}}{2(m+1)b_{w}} $ &  $\frac{1+(1+k\rho-\rho-2k)w_{G}}{2(m+1)b_{w}}$ \\
 $t_{1}<t<t_{2}$  & $\frac{1-(1+k\rho-\rho)w_{G}}{2}$ & $ \frac{(mw_{G}k+1)\rho-(m+1)(\bar{c}-t)}{2 b_{w}\rho} $  &$ \frac{1+(-1+k\rho-\rho)w_{G}}{2(m+1)b_{w}} $ & $\frac{1+(1+k\rho-\rho-2k)w_{G}}{2(m+1)b_{w}}-\frac{Q_{s}^{E}}{m+1}$  \\
 $t_{2}\leq t<\bar{c}$ & $\frac{1-w_{G}}{2}$ & $\frac{\rho+t-\bar{c}}{2 b_{w}\rho}$ &$ \frac{1-w_{G}}{2(m+1)b_{w}} $ & 0  \\  $0<t\leq t_{1}$ & $\frac{1-(1+k\rho-\rho)w_{G}}{2}$ & 0 &$ \frac{1+(-1+k\rho-\rho)w_{G}}{2(m+1)b_{w}} $ &  $\frac{1+(1+k\rho-\rho-2k)w_{G}}{2(m+1)b_{w}}$ \\
 $t_{1}<t<t_{2}$  & $\frac{1-(1+k\rho-\rho)w_{G}}{2}$ & $ \frac{(mw_{G}k+1)\rho-(m+1)(\bar{c}-t)}{2 b_{w}\rho} $  &$ \frac{1+(-1+k\rho-\rho)w_{G}}{2(m+1)b_{w}} $ & $\frac{1+(1+k\rho-\rho-2k)w_{G}}{2(m+1)b_{w}}-\frac{Q_{s}^{E}}{m+1}$  \\
 $t_{2}\leq t<\bar{c}$ & $\frac{1-w_{G}}{2}$ & $\frac{\rho+t-\bar{c}}{2 b_{w}\rho}$ &$ \frac{1-w_{G}}{2(m+1)b_{w}} $ & 0  \\  $0<t\leq t_{1}$ & $\frac{1-(1+k\rho-\rho)w_{G}}{2}$ & 0 &$ \frac{1+(-1+k\rho-\rho)w_{G}}{2(m+1)b_{w}} $ &  $\frac{1+(1+k\rho-\rho-2k)w_{G}}{2(m+1)b_{w}}$ \\
 $t_{1}<t<t_{2}$  & $\frac{1-(1+k\rho-\rho)w_{G}}{2}$ & $ \frac{(mw_{G}k+1)\rho-(m+1)(\bar{c}-t)}{2 b_{w}\rho} $  &$ \frac{1+(-1+k\rho-\rho)w_{G}}{2(m+1)b_{w}} $ & $\frac{1+(1+k\rho-\rho-2k)w_{G}}{2(m+1)b_{w}}-\frac{Q_{s}^{E}}{m+1}$  \\
 $t_{2}\leq t<\bar{c}$ & $\frac{1-w_{G}}{2}$ & $\frac{\rho+t-\bar{c}}{2 b_{w}\rho}$ &$ \frac{1-w_{G}}{2(m+1)b_{w}} $ & 0  \\  $0<t\leq t_{1}$ & $\frac{1-(1+k\rho-\rho)w_{G}}{2}$ & 0 &$ \frac{1+(-1+k\rho-\rho)w_{G}}{2(m+1)b_{w}} $ &  $\frac{1+(1+k\rho-\rho-2k)w_{G}}{2(m+1)b_{w}}$ \\
 $t_{1}<t<t_{2}$  & $\frac{1-(1+k\rho-\rho)w_{G}}{2}$ & $ \frac{(mw_{G}k+1)\rho-(m+1)(\bar{c}-t)}{2 b_{w}\rho} $  &$ \frac{1+(-1+k\rho-\rho)w_{G}}{2(m+1)b_{w}} $ & $\frac{1+(1+k\rho-\rho-2k)w_{G}}{2(m+1)b_{w}}-\frac{Q_{s}^{E}}{m+1}$  \\
 $t_{2}\leq t<\bar{c}$ & $\frac{1-w_{G}}{2}$ & $\frac{\rho+t-\bar{c}}{2 b_{w}\rho}$ &$ \frac{1-w_{G}}{2(m+1)b_{w}} $ & 0  \\  $0<t\leq t_{1}$ & $\frac{1-(1+k\rho-\rho)w_{G}}{2}$ & 0 &$ \frac{1+(-1+k\rho-\rho)w_{G}}{2(m+1)b_{w}} $ &  $\frac{1+(1+k\rho-\rho-2k)w_{G}}{2(m+1)b_{w}}$ \\
 $t_{1}<t<t_{2}$  & $\frac{1-(1+k\rho-\rho)w_{G}}{2}$ & $ \frac{(mw_{G}k+1)\rho-(m+1)(\bar{c}-t)}{2 b_{w}\rho} $  &$ \frac{1+(-1+k\rho-\rho)w_{G}}{2(m+1)b_{w}} $ & $\frac{1+(1+k\rho-\rho-2k)w_{G}}{2(m+1)b_{w}}-\frac{Q_{s}^{E}}{m+1}$  \\
 $t_{2}\leq t<\bar{c}$ & $\frac{1-w_{G}}{2}$ & $\frac{\rho+t-\bar{c}}{2 b_{w}\rho}$ &$ \frac{1-w_{G}}{2(m+1)b_{w}} $ & 0  \\
\bottomrule
\end{tabular}%}
\end{table}



3)当竖向项目较少,而横向栏目过多,或宽度过大时,可以将表格转换成上下重叠的2段或3段, 段与段之间用双细线隔开,每段左方的横向项目应重复编排.

\begin{table}[H]
\centering
\caption{政策调整前后深度虚值期权对未来波动的影响对比}
\label{tab6}
\tabcolsep 9pt
    \begin{tabular}{rlllllr}
    \toprule
          & \multicolumn{6}{c}{政策调整前(2015--2017)} \\
          \midrule
   & \multicolumn{1}{c}{$\Phi_p$} & \multicolumn{1}{c}{$\Phi_c$}& \multicolumn{1}{c}{$\beta_d$} & \multicolumn{1}{c}{$\beta_w$}& \multicolumn{1}{c}{$\beta_m$} & \multicolumn{1}{c}{Adj-$R^2$}\\
                \midrule
    \multicolumn{1}{l}{$h=1$} &       &       & 0.394*** & 0.362*** & 0.181*** & \multicolumn{1}{l}{0.771} \\
          &       &       & (8.485) & (5.386) & (3.726) &  \\
          & 0.023 & $-$0.014 & 0.391*** & 0.359*** & 0.188*** & \multicolumn{1}{l}{0.771} \\
          & (0.787) & ($-$0.590) & (8.376) & (5.302) & (3.836) &  \\
    \multicolumn{1}{l}{$h=5$} &       &       & 0.230*** & 0.311*** & 0.326*** & \multicolumn{1}{l}{0.640} \\
          &       &       & (3.945) & (3.684) & (5.350) &  \\
          & 0.104*** & $-$0.045 & 0.228*** & 0.287*** & 0.351*** & \multicolumn{1}{l}{0.648} \\
          & (2.818) & ($-$1.481) & (3.944) & (3.412) & (5.772) &  \\
                \hline\hline
          & \multicolumn{6}{c}{政策调整后(2018$-$2021)} \\
                \midrule
  & \multicolumn{1}{c}{$\Phi_p$} & \multicolumn{1}{c}{$\Phi_c$}& \multicolumn{1}{c}{$\beta_d$} & \multicolumn{1}{c}{$\beta_w$}& \multicolumn{1}{c}{$\beta_m$} & \multicolumn{1}{c}{Adj-$R^2$}\\
                \midrule
    \multicolumn{1}{l}{$h=1$} &       &       & 0.279*** & 0.526*** & 0.035 & \multicolumn{1}{l}{0.502} \\
          &       &       & (6.731) & (8.570) & (0.673) &  \\
          & 0.056* & $-$0.049* & 0.268*** & 0.550*** & 0.038 & \multicolumn{1}{l}{0.504} \\
          & (1.724) & ($-$1.658) & (6.438) & (8.821) & (0.736) &  \\
    \multicolumn{1}{l}{$h=5$} &       &       & 0.268*** & 0.325*** & 0.044 & \multicolumn{1}{l}{0.290} \\
          &       &       & (5.418) & (4.446) & (0.708) &  \\
          & 0.097** & $-$0.190*** & 0.244*** & 0.373*** & 0.054 & \multicolumn{1}{l}{0.319} \\
          & (2.532) & ($-$5.445) & (4.992) & (5.116) & (0.889) &  \\
    \bottomrule
\end{tabular}
\end{table}

4)当横向栏目较少,而竖向项目过多时, 可以把它拦腰截断后平行地转换成2列或3列的形式,这里的列与列之间用双细线隔开.每列的栏头应重复排出.


\begin{table}[H]

\centering
\caption{建立模型所用变量}%\tabcolsep 2pt
\label{tab1}
	\begin{tabular}{cc||cc}\toprule
			序号  &       变量名称 &       序号 &                 变量名称 \\\hhline{--||--}
			$x_1$  &       搅拌速率   &   $x_7$  &      碱流加速率      \\
			$x_2$  &       通风速率    &   $x_8$  &      酸流加速率      \\
			$x_3$  &       补料温度     &   $x_9$  &       温度\\
			$x_4$  &       底物流加速率 &  ${x_{10}}$  &       冷水流速\\
			$x_5$  &        溶氧浓度    &  ${x_{11}}$  &       菌体浓度\\
			$x_6$  &        PH值              &  ${x_{12}}$  &       产物浓度\\
			\bottomrule
	\end{tabular}
\end{table}




5)当某表格1页排不下时可在下一页排续表
\begin{table}[H]
\centering
\caption*{表4\ \ (续)}%\resizebox{\textwidth}{!}{
%\tabcolsep 2pt%\renewcommand{\arraystretch}{0.98}
\begin{tabular}{lcccc} \toprule
     \multirow{2}{*}{条件} & \multirow{2}{*}{单位交易费用$(\alpha^{E})$} &  \multirow{2}{*}{储备数量$(Q_{s}^{E})$} & \multicolumn{2}{c}{批发商$i$的订货数量}  \\\cmidrule(r){4-5}
     &&& 正常情形$(q_{iG}^{E})$ & 突发事件情形$(q_{iL}^{E})$ \\

\hline
  $0<t\leq t_{1}$ & $\frac{1-(1+k\rho-\rho)w_{G}}{2}$ & 0 &$ \frac{1+(-1+k\rho-\rho)w_{G}}{2(m+1)b_{w}} $ &  $\frac{1+(1+k\rho-\rho-2k)w_{G}}{2(m+1)b_{w}}$ \\
 $t_{1}<t<t_{2}$  & $\frac{1-(1+k\rho-\rho)w_{G}}{2}$ & $ \frac{(mw_{G}k+1)\rho-(m+1)(\bar{c}-t)}{2 b_{w}\rho} $  &$ \frac{1+(-1+k\rho-\rho)w_{G}}{2(m+1)b_{w}} $ & $\frac{1+(1+k\rho-\rho-2k)w_{G}}{2(m+1)b_{w}}-\frac{Q_{s}^{E}}{m+1}$  \\
 $t_{2}\leq t<\bar{c}$ & $\frac{1-w_{G}}{2}$ & $\frac{\rho+t-\bar{c}}{2 b_{w}\rho}$ &$ \frac{1-w_{G}}{2(m+1)b_{w}} $ & 0  \\
\bottomrule
\end{tabular}
\end{table}


6) 表格转置(不推荐)
\begin{landscape}
\begin{table}[h]
\centering
\caption{企业异质性: 企业基本特征}
\label{tab11}
\tabcolsep 9pt
    \begin{tabular}{p{2.5cm}<{\raggedright} ccccccccccc}
    \toprule
    因变量:   & \multicolumn{11}{c}{lnwait} \\
\cmidrule{2-12}          & \multicolumn{3}{c}{盈利能力} &       & \multicolumn{3}{c}{负债比例} &       & \multicolumn{3}{c}{公司规模} \\
\cmidrule{2-4}\cmidrule{6-8}\cmidrule{10-12}    分样本:   & 低盈利   & 高盈利   & 全样本   &       & 高负债   & 低负债   & 全样本   &       & 大规模   & 小规模   & 全样本 \\
\cmidrule{2-4}\cmidrule{6-8}\cmidrule{10-12}          & (1)   & (2)   & (3)   &       & (4)   & (5)   & (6)   &       & (7)   & (8)   & (9) \\
    \midrule
          &       &       &       &       &       &       &       &       &       &       &  \\
    VCgov & -0.218** & -0.273** & -0.343*** &       & -0.226** & -0.234** & -0.329*** &       & -0.459*** & -0.144* & -0.177** \\
          & (0.111) & (0.114) & (0.119) &       & (0.108) & (0.106) & (0.110) &       & (0.135) & (0.084) & (0.090) \\
    lprofit &       &       & -0.117 &       &       &       & \multicolumn{1}{l}{} &       &       &       & \multicolumn{1}{l}{} \\
          &       &       & (0.080) &       &       &       & \multicolumn{1}{l}{} &       &       &       & \multicolumn{1}{l}{} \\
    VCgov * lprofit &       &       & 0.080 &       &       &       & \multicolumn{1}{l}{} &       &       &       & \multicolumn{1}{l}{} \\
          &       &       & (0.158) &       &       &       & \multicolumn{1}{l}{} &       &       &       & \multicolumn{1}{l}{} \\
   hdebt &       &       &       &       &       &       & -0.168** &       &       &       & \multicolumn{1}{l}{} \\
          &       &       &       &       &       &       & (0.066) &       &       &       & \multicolumn{1}{l}{} \\
  VCgov * hdebt &       &       &       &       &       &       & 0.050 &       &       &       & \multicolumn{1}{l}{} \\
          &       &       &       &       &       &       & (0.155) &       &       &       & \multicolumn{1}{l}{} \\
    hsize &       &       &       &       &       &       & \multicolumn{1}{l}{} &       &       &       & 0.082 \\
          &       &       &       &       &       &       & \multicolumn{1}{l}{} &       &       &       & (0.085) \\
   VCgov * hsize &       &       &       &       &       &       & \multicolumn{1}{l}{} &       &       &       & -0.285* \\
          &       &       &       &       &       &       & \multicolumn{1}{l}{} &       &       &       & (0.164) \\
          &       &       &       &       &       &       &       &       &       &       &  \\
    基本控制变量 & 控制    & 控制    & 控制    &       & 控制    & 控制    & 控制    &       & 控制    & 控制    & 控制 \\
    行业、申请月份和省份固定效应 & 控制    & 控制    & 控制    &       & 控制    & 控制    & 控制    &       & 控制    & 控制    & 控制 \\
    观测值个数 & 610   & 610   & 1,220 &       & 704   & 516   & 1,220 &       & 610   & 610   & 1,220 \\
    调整后$R^2$ & 0.755 & 0.706 & 0.685 &       & 0.783 & 0.635 & 0.686 &       & 0.551 & 0.843 & 0.685 \\
    \bottomrule
    \end{tabular}%
\end{table}
\end{landscape}



\subsection{参考文献}


1) 参考文献在正文中著录格式的基本要求
为反映论文的学术水平和创新程度, 请主要引用最近3年以内发表的文献(最好引用期刊类文献). 作者可在本刊网站: http://www.sysengi.com 免费下载查阅创刊以来的所有论文全文.

2)参考文献列表必须按正文中引用的先后顺序排序, 所有被引用的文献须在正文中标明引用位置: 不做句子成份的文献以上角标形式标引, 做句子成份的文献, 则与正文平排.

示例:

``本文利用粒子群算法[1]的优点, 结合文献[2]所构造的算法, 提出了一种改进的优化算法."

3) 文献类型标志

C: 会议录; D: 学位论文; G: 汇编; J: 期刊; M: 普通图书; N: 报纸; P: 专利; R: 报告; S: 标准; EB: 电子公告.


{\bf 重点:}

4)参考文献列表中的每条文献必须在正文中被引用, 文献列表按正文中被引用的顺序排列. 应采用顺序编码制, 不要与著者出版年制混合, 即不要出现孙会霞(2021)~\cite{ref1}.应为孙会霞等~\cite{ref1}.

5) 文献引用要与参考文献一一对应.

6) 非英文文献者均需译成英文附在其后.

7)  文献起始页码用\verb|--| 如: 2022, 33(1): 189--190;
中文破折号用\verb|------|; 英文破折号用\verb|---|.

如: 孙会霞, 苏峻, 何佳. 股票供给控制、~需求曲线与股价的反应------基于创业板的经验数据\\
 Control on supply of shares, demand curves and the reaction of stock price --- Based on the data of Chinext.

8) 若文献为网络首发或文献的年卷期不全, 请加上文章的doi: $\cdots$$\cdots$.

文献示例如下:

\begin{thebibliography}{99}


%1.期刊中的析出文献
%格式: 析出文献主要责任者.析出文献题名 [文献类型标志]. 连续出版物题名:其他题名信息, 年, 卷(期): 起页-止页. (为方便EI检索, 中文文献如有英文文题时请同时给出对应的英文表述, 不要自行翻译)
%
%提示:
%
%a. 期刊名称应采用全称, 不采用缩写形式表示.
%
%b. 出版年、卷、期、页码须齐全, 正确著录格式如下: 年, 卷(期): 页. 例如: 2005, 10(2): 15-20.
%
%c. 文献题名中首词首字母、缩略语、专有名词(如Monte Carlo)首字母外, 均小写; 期刊名各实词首字母大写.
%
%示例:
\bibitem{ref1} 孙会霞, 苏峻, 何佳. 股票供给控制、~需求曲线与股价的反应------基于创业板的经验数据[J]. 系统工程理论与实践, 2013, 33(1): 1--11. \\
Sun H X, Su J, He J. Control on supply of shares, demand curves and the reaction of stock price --- Based on the data of Chinext[J]. Systems Engineering --- Theory \& Practice, 2013, 33(1): 1--11.


%2.专著、会议录、论文集中的析出文献
%
%格式: 析出文献主要责任者. 析出文献题名[文献类型标志]//原文献主要责任者.专著、会议录、论文集题名, 出版地: 出版者, 出版年: 起页-止页.
%
%提示:
%a. 会议论文一律按专著中的析出文献的格式著录, 特别是要正确著录会议文集的名称和出版年、起止页码.
%b. 会议文集名称应采用全称, 不可以用简写和缩写形式.
%
%示例:
\bibitem{ref2} 席酉民, 王亚刚. 和谐社会秩序形成机制的系统分析: 和谐管理理论的启示和价值[C]// 中国系统工程学会第十四届学术年会论文集, 2006: 18--29.
%
%
%3.普通图书
%格式: 主要责任者. 题名[M]. 版本项. 出版地: 出版者, 出版年: 起页-止页.
%
%提示: 如果是第1版版本项可以省略, 第2版中文著录格式为: 2版, 英文著录格式为: 2nd, ed.
%
%示例:
\bibitem{ref3} 徐玖平, 胡知能. 经济管理的动态理论: 稳定与失衡[M]. 北京: 科学出版社, 2004.\\
Xu J P, Hu Z N. The development theories of economy management: Stability and imbalance[M]. Beijing: Science Press, 2004.


%4. 学位论文
%
%格式: 主要责任者. 题名[D]. 保存地: 保存单位, 年份.
%
%示例:
\bibitem{ref4} 刘舒文. 会计稳健的检验及其影响因素和经济后果的分析[D]. 长沙: 湖南大学, 2007.\\
Liu S W. The test of accounting conservatism and its affecting factors analysis and economic
consequence analysis[D]. Changsha: Hunan University, 2007.


%5. 科技报告
%
%格式: 主要责任者. 题名[R]. 出版地: 出版者, 出版年.
%
%示例:
\bibitem{ref5} Chen L. Corporate yield spreads and bond liquidity[R]. East Lansing: Michigan State University, 2005.


%6. 专 利
%
%格式: 专利所有者. 专利题名[P]. 专利国别: 专利号, 发布日期.
%
%示例:
\bibitem{ref6} 姜锡洲. 一种温热外敷药制备方案[P]. 中国: 881056078, 1983--08--12.
\\Jiang X Z. A warm external compress preparation scheme[P]. China:  881056078, 1983--08--12.

%7.国际标准和国家标准
%
%格式: 标准代号. 标准名称[S]. 出版地: 出版者, 出版年.
%
%示例:
\bibitem{ref7} GB/T 16159-1996. 汉语拼音正词法基本规则[S]. 北京: 中国标准出版社, 1996.
GB/T 16159-1996. Basic rules of Hanyu Pinyin orthography[S]. Beijing: Standards Press of China, 1996.


%8. 电子文献
%
%格式: 主要责任者. 文题[EB/OL].[引用日期]. 获取和访问路径.
%
%示例:
\bibitem{ref8}萧钰.出版业信息化迈入快车道[EB/OL]. [2002--04--15].\\
http://www.booktide.com/news/20011219/200112190019.html.\\
Xiao Y. The informatization of the publishing industry has entered the fast lane[EB/OL]. [2002--04--15].
http://www.booktide.com/news/20011219/200112190019.html.\\

%提示: 方括号[]中的年月日为文献被引用的时间, 非网站上的更新日期;此处提到的电子文献指网上文献, 光盘类文献的标志代码为[EB/CD].

\end{thebibliography}



\clearpage
 
\appendix



\section*{附录}
  %文章分为正文和附录两部分, 请作者将附录直接附在正文的最后面, 文章发表时纸质期刊仅刊发正文部分; 本刊官网、~中国知网等电子版将刊发全文(正文+附录). 在不影响可读性情况下, 文章背景介绍、~问题描述、~分析过程等尽量言简意赅, 请尽可能将证明过程、~算法步骤(/步骤图)、~数据分析过程、~源数据、~成熟的概念或模型等放入附录.
\pagenumbering{Roman}

\setcounter{page}{1}

\section{附录1}







\newpage%另起一页


  \section{附录2}

\end{document}
